\documentclass[journal]{IEEEtran}

% Packages
\usepackage{cite}
\usepackage{graphicx}
\usepackage{amsmath}
\usepackage{algorithmic}
\usepackage{array}
\usepackage{url}
\usepackage{hyperref}
\usepackage{booktabs}
\usepackage{multirow}

\begin{document}

% ==========================================================================
\title{Evaluating LLM-Based Code and Test Generation on HumanEval-Java:\\
A Comparative Study of ChatGPT and Gemini}

\author{\IEEEauthorblockN{Oğuz Eren Kacar, <Student 2>, <Student 3>}\\
\IEEEauthorblockA{Department of Computer Engineering\\
Istanbul Technical University\\
Istanbul, Turkey\\
\{150200018, <ID2>, <ID3>\}@itu.edu.tr}}

\maketitle

% ==========================================================================
\begin{abstract}
The rapid adoption of large language models (LLMs) in software development
makes it essential to rigorously validate the code they produce. This study
evaluates two state-of-the-art LLMs OpenAI ChatGPT (GPT-4o) and Google
Gemini 2.5 Pro on a carefully selected 30-problem subset of the
HumanEval-Java benchmark, spanning easy, moderate, and hard difficulty levels.
Using a semi-agentic workflow, we prompt each model to generate Java solutions,
transliterate the provided HumanEval base tests into JUnit~5/6, and then
systematically improve those tests through equivalence class partitioning and
boundary value analysis. Correctness is measured by pass/fail outcomes;
structural quality is assessed via JaCoCo branch coverage. ChatGPT (GPT-4o)
produced correct implementations for all 30 problems, while Gemini 2.5 Pro
succeeded on only 23 of 30, with failures attributable to off-by-one errors,
structural misinterpretations, and divergent edge-case contracts. Applying our
improved EC/BV-guided test suite achieved 100\% branch coverage across every
ChatGPT-generated class that contains conditional branches, demonstrating that
systematic black-box analysis substantially strengthens LLM-generated test
suites beyond the baseline.
\end{abstract}

\begin{IEEEkeywords}
LLM, code generation, test generation, HumanEval, JUnit, branch coverage,
equivalence class partitioning, boundary value analysis.
\end{IEEEkeywords}

% ==========================================================================
\section{Introduction}
\label{sec:intro}

Large Language Models (LLMs) have revolutionized software development by enabling developers to generate code from natural language descriptions. Tools like ChatGPT and Gemini can automatically produce executable code, significantly accelerating prototyping and reducing development time~\cite{yuan2025testctrl}. However, the correctness and quality of LLM-generated code remain critical concerns that necessitate rigorous testing.

The HumanEval-X dataset~\cite{peng2024humanevalxl} was introduced as a multilingual extension of the popular HumanEval benchmark. It provides a standardised set of programming problems in multiple natural languages and programming languages, allowing researchers to assess the cross-lingual generalisation of code generation models. We focus on the Java subset of this benchmark, which poses a diverse set of challenges from simple string manipulations to complex algorithmic tasks.

In this work, we address two main research questions: (1) How do the latest versions of ChatGPT and Gemini perform in terms of code correctness on 30 HumanEval-Java problems? (2) What are the common failure patterns of the weaker LLM, and how effective is a semi‑agentic, coverage-guided test improvement loop? Our primary contribution is a systematic, side‑by‑side comparison that not only measures pass/fail rates but also analyses line and branch coverage, test quality, and the impact of manual prompt refinement. To the best of our knowledge, this is one of the first studies to directly compare these two leading commercial LLMs on a Java‑specific benchmark using a structured testing methodology.


% ==========================================================================
% This file is INCLUDED from main.tex via % This file is INCLUDED from main.tex via % This file is INCLUDED from main.tex via \input{literature_review}.
% TODO (Member B): Replace the placeholder bullets with real paraphrased
% summaries of your 5 papers. Do NOT use LLM tools for this section; the
% assignment explicitly forbids it. Paraphrase in your own voice.

% Suggested structure for each paper (3--4 sentences):
%   (1) What problem did they tackle?
%   (2) What method did they use?
%   (3) What was the headline result?
%   (4) How does it relate to our study?

Research on LLM-based code and test generation has accelerated since late
2022. We summarise five representative studies published between 2023 and
2026 that we found through Google Scholar and the IEEE Xplore digital
library.

\paragraph*{Paper 1 [TODO \cite{paper1}]}
TODO (Member B): 3--4 sentence paraphrase.

\paragraph*{Paper 2 [TODO \cite{paper2}]}
TODO.

\paragraph*{Paper 3 [TODO \cite{paper3}]}
TODO.

\paragraph*{Paper 4 [TODO \cite{paper4}]}
TODO.

\paragraph*{Paper 5 [TODO \cite{paper5}]}
TODO.

\medskip
Compared with these works, our study is narrower --- 30 problems in Java,
two LLMs, no automated feedback loop --- but complements them by providing
paired base/improved test suites and a reproducible per-problem interaction
log.
.
% TODO (Member B): Replace the placeholder bullets with real paraphrased
% summaries of your 5 papers. Do NOT use LLM tools for this section; the
% assignment explicitly forbids it. Paraphrase in your own voice.

% Suggested structure for each paper (3--4 sentences):
%   (1) What problem did they tackle?
%   (2) What method did they use?
%   (3) What was the headline result?
%   (4) How does it relate to our study?

Research on LLM-based code and test generation has accelerated since late
2022. We summarise five representative studies published between 2023 and
2026 that we found through Google Scholar and the IEEE Xplore digital
library.

\paragraph*{Paper 1 [TODO \cite{paper1}]}
TODO (Member B): 3--4 sentence paraphrase.

\paragraph*{Paper 2 [TODO \cite{paper2}]}
TODO.

\paragraph*{Paper 3 [TODO \cite{paper3}]}
TODO.

\paragraph*{Paper 4 [TODO \cite{paper4}]}
TODO.

\paragraph*{Paper 5 [TODO \cite{paper5}]}
TODO.

\medskip
Compared with these works, our study is narrower --- 30 problems in Java,
two LLMs, no automated feedback loop --- but complements them by providing
paired base/improved test suites and a reproducible per-problem interaction
log.
.
% TODO (Member B): Replace the placeholder bullets with real paraphrased
% summaries of your 5 papers. Do NOT use LLM tools for this section; the
% assignment explicitly forbids it. Paraphrase in your own voice.

% Suggested structure for each paper (3--4 sentences):
%   (1) What problem did they tackle?
%   (2) What method did they use?
%   (3) What was the headline result?
%   (4) How does it relate to our study?

Research on LLM-based code and test generation has accelerated since late
2022. We summarise five representative studies published between 2023 and
2026 that we found through Google Scholar and the IEEE Xplore digital
library.

\paragraph*{Paper 1 [TODO \cite{paper1}]}
TODO (Member B): 3--4 sentence paraphrase.

\paragraph*{Paper 2 [TODO \cite{paper2}]}
TODO.

\paragraph*{Paper 3 [TODO \cite{paper3}]}
TODO.

\paragraph*{Paper 4 [TODO \cite{paper4}]}
TODO.

\paragraph*{Paper 5 [TODO \cite{paper5}]}
TODO.

\medskip
Compared with these works, our study is narrower --- 30 problems in Java,
two LLMs, no automated feedback loop --- but complements them by providing
paired base/improved test suites and a reproducible per-problem interaction
log.


% ==========================================================================
\section{Methodology}
\label{sec:method}

\subsection{LLM Selection Rationale}
% TODO (Member B). Key points (from WORK_DIVISION.md):
%   - ChatGPT (GPT-4o) is the industry baseline.
%   - Gemini 2.5 Pro is Google's flagship alternative; gives vendor diversity.
%   - Both are publicly accessible without paid research credentials.
%   - Both expose a plain chat interface --- fits the semi-agentic workflow.

\subsection{Approach: Semi-Agentic vs Agentic}
% TODO (Member B). Justify the semi-agentic choice: human-in-the-loop keeps
% prompts identical across both LLMs, which is necessary for a fair comparison.
% Limitations: doesn't exercise iterative self-correction; slower.

\subsection{Problem Selection}
We selected 30 problems from the 164-problem HumanEval-Java
dataset~\cite{humaneval-x} covering three difficulty categories:

\begin{itemize}
\item \textbf{Easy (10)}: HE-000, 007, 009, 012, 021, 023, 024, 027, 028, 042
\item \textbf{Moderate (10)}: HE-003, 011, 014, 015, 017, 018, 025, 031, 034, 035
\item \textbf{Hard (10)}: HE-010, 040, 047, 055, 072, 083, 109, 126, 132, 139
\end{itemize}

Difficulty was assigned by inspection of the canonical solution complexity:
easy problems are one-pass linear algorithms on primitives or short strings;
moderate problems involve multi-branch conditionals or nested loops; hard
problems require non-trivial invariants, mathematical reasoning, or
multi-case boundary handling.

\subsection{Code Generation}
For each problem we issued the identical prompt (see
\texttt{logs/<llm>/HE-XXX\_codegen.md}) to both LLMs. The response was pasted
into the source tree without modification, per the assignment's constraint
at this step.

\subsection{Test Conversion and Improvement}
HumanEval's base tests are written as a single \texttt{main} that assembles
a list of boolean assertions. We transliterated each into JUnit 5/6 using
\texttt{assertTrue} with one assertion per original boolean (see
\texttt{*BaseTest.java}). Expected values were not changed.

For the improved suite we performed equivalence class partitioning and
boundary value analysis per problem (see
Section~\ref{sec:ec-bv-tables} and \texttt{report/equivalence\_classes.md})
and wrote targeted \texttt{@Test} methods, one per partition or boundary.

\subsection{Coverage Measurement}
JaCoCo 0.8.12 reports branch coverage; results are summarised in
Section~\ref{sec:results}.

% ==========================================================================
\section{Results}
\label{sec:results}

\subsection{Pass Rate Per LLM}
\begin{table}[h]
\centering
\caption{Base and Improved Test Pass Rate}
\label{tab:pass-rate}
\begin{tabular}{lccc}
\toprule
LLM & Base tests & Improved tests & Total pass \\
\midrule
ChatGPT (GPT-4o) & 30/30 & 30/30 & 60/60 \\
Gemini 2.5 Pro   & 24/30 & 23/30 & 47/60 \\
\bottomrule
\end{tabular}
\end{table}

% Actual numbers from our run: ChatGPT 30/30 & 30/30, Gemini: 24/30 base
% (HE-000, 040, 072, 083, 132, 139 fail), 23/30 improved (those six plus
% HE-018 due to divergent empty-substring contract).

\subsection{Branch Coverage}
Table~\ref{tab:coverage} summarises branch coverage from
\texttt{jacocoTestReport}. On all ChatGPT classes that contain conditional
branches, the improved suite reaches 100\% branch coverage. Classes with
``N/A'' have no conditionals in the canonical solution (e.g.
\texttt{strlen} is a single \texttt{return} statement).

% TODO (Member B): Import tabular data from
% report/coverage/coverage_summary.md.

\subsection{Problem-Level Outcomes}
% TODO (Member B): Summarise which 7 Gemini implementations failed and why,
% drawing from logs/gemini/HE-*_codegen.md. See the WORK_DIVISION.md
% description of each bug class.

% ==========================================================================
\section{Discussion}
\label{sec:discuss}

\subsection{Nature of Gemini Deviations}
% TODO (Member B): 1 paragraph per bug class (6 bug classes observed):
% 1. Missing palindrome check (HE-072 willItFly)
% 2. Off-by-one in loop bound (HE-040 triplesSumToZero, HE-139 specialFactorial)
% 3. Incorrect closed-form (HE-083 startsOneEnds)
% 4. Structural misinterpretation (HE-132 isNested as bracket-count)
% 5. Stream filter masking equality (HE-000 hasCloseElements)
% 6. Divergent edge-case contract (HE-018 empty substring)

\subsection{Effectiveness of EC/BV Partitioning}
% TODO (Member B): Observations:
%   - HE-132 base test didn't cover the "siblings" partition; improved did.
%   - HE-018 edge-case contract divergence only surfaced in the improved suite.
% EC/BV partitioning was therefore strictly stronger than the HumanEval base
% tests for at least these two problems.

\subsection{Limitations}
% TODO (Member B): Sample size = 30 problems. Only two LLMs. Semi-agentic
% (no iterative correction measured). Both LLMs accessed via chat UI, not
% API, so temperature/seed not controlled.

\subsection{Refactoring Iteration}
% TODO (Member B): Describe Member C's refactoring log. For each of the 7
% Gemini failures, how many prompt rounds were needed to converge? What
% prompt phrasing worked best?

% ==========================================================================
\section{Conclusion}
\label{sec:conclude}
% TODO (Member B): 1 paragraph, summarising findings and pointing to Phase 2.

% ==========================================================================
\section*{Acknowledgment}
\label{sec:ack}
% TODO (Member C): Fill in the table with final names/IDs and concrete
% per-person tasks actually performed.

The authors thank the instructor Sultan \c{C}o\u{g}ay for organising the
course and for the detailed problem specification. We acknowledge the
public availability of the HumanEval-X dataset, which made this study
possible. The repository for this project is available at:
\begin{center}
\url{https://github.com/<username>/blg475e-project-2026}
\end{center}

\subsection*{Work Distribution}

\begin{tabular}{p{0.18\linewidth} p{0.75\linewidth}}
\toprule
Member & Contribution \\
\midrule
<Student 1> (Mert Ayd\i n) &
Pipeline and build system; problem selection script; automated generation of
ChatGPT and Gemini sources, base tests, and improved tests; LLM interaction
log templating; JaCoCo configuration; coverage summary script. \\
\addlinespace
<Student 2> &
Literature review and bibliographic search (five papers, 2023--2026);
drafting of the IEEE report: Abstract, Introduction, Related Work,
Methodology, Results, Discussion, and Conclusion; acknowledgments layout. \\
\addlinespace
<Student 3> &
Manual review of equivalence-class partitions and boundary-value tests;
refactoring iterations with Gemini for the seven divergent implementations;
collection of JaCoCo coverage screenshots; final quality-assurance pass on
source-file headers and commit history. \\
\bottomrule
\end{tabular}


% ==========================================================================
\begin{thebibliography}{00}
\bibliography{references}
\bibliographystyle{IEEEtran}
\end{thebibliography}

\end{document}
